\documentclass[12pt]{article}
\usepackage[utf8]{inputenc}
\usepackage[T1]{fontenc}
\usepackage{hyperref}
\title{Yareba Language Resources for K-12 Curriculum Development}
\author{Compiled by Claude Sonnet 5.5 (claude-sonnet-5-5) for Andr\'as Kornai}
\date{}
\begin{document}
\maketitle

\section{Introduction}
Yareba (Middle Musa) is a Yareban, Trans--New Guinea language spoken in the Middle Musa valley, Oro Province, Papua New Guinea. Estimates range from roughly 800--1,200 speakers to 2,600 for the people group (Joshua Project). Glottolog lists it as not endangered. Wikipedia reports 15 consonant phonemes and 5 monophthongs.

\section{Primary Educational Resources}
No primers, textbooks or school programmes were found in the searches.

\section{Linguistic and Reference Materials}
\begin{itemize}
\item Harry and Natalia Weimer, \emph{Yareba} dictionary (SIL, Dictionaries of Papua New Guinea series; catalogue dates of about 1974 and 1977). Only catalogue records were found, not full text. Glottolog reference record: \url{https://glottolog.org/resource/reference/id/302898}. Library record: \url{https://matdisc.unitech.ac.pg/Record/9344}
\item Wikipedia on Yareba: \url{https://en.wikipedia.org/wiki/Yareba_language}
\item Wikipedia on the Yareban languages: \url{https://en.wikipedia.org/wiki/Yareban_languages}
\item ASJP word list for Yareba: \url{https://asjp.clld.org/languages/YAREBA}
\item Max Planck numeral page for Yareba: \url{https://lingweb.eva.mpg.de/channumerals/Yareba.htm}
\end{itemize}

\section{Religious and Community Resources}
\begin{itemize}
\item Joshua Project profile of the Yareba (Middle Musa): \url{https://joshuaproject.net/people_groups/15994}
\item Joshua Project language page for yrb: \url{https://joshuaproject.net/index.php/languages/yrb}
\end{itemize}
The search summary says the New Testament and some Scripture portions exist in Yareba. No direct scripture link appeared in the results.

\section{Digital and Audio Resources}
No digital or audio resources were identified in the searches.

\section{Curriculum Development Context}
Yareba is reportedly still the main village language, which is a good basis for mother-tongue education. No school materials were found. The Weimer dictionary and the Scripture translation are the likely starting points, and their availability and rights should be checked with SIL.
\end{document}